\documentclass[a4paper,12pt]{article}

\usepackage{polyglossia} 
\setmainlanguage{russian}
\usepackage{fontspec}
\setmainfont{Times New Roman}
\usepackage{geometry}
\usepackage{setspace}
\usepackage{amsmath}
\usepackage{graphicx} 
\usepackage{listings}
\usepackage{xcolor}
\setmonofont{Courier New}
\newfontfamily{\cyrillicfonttt}{Courier New}


\geometry{top=20mm, bottom=20mm, left=25mm, right=15mm}

\lstset{
    language=Python,
    basicstyle=\ttfamily\scriptsize,
    keywordstyle=\color{blue},
    stringstyle=\color{red},
    commentstyle=\color{green!60!black},
    numbers=left,
    numberstyle=\tiny\color{gray},
    numbersep=5pt,
    backgroundcolor=\color{gray!5},
    breaklines=true,
    frame=single
}

\begin{document}

\noindent
\begin{minipage}[c]{0.22\textwidth}
    \centering
    
\includegraphics[width=\textwidth]{../bmstu_logo.jpg} 
\end{minipage}%
\hfill
\begin{minipage}[c]{0.75\textwidth}
    \centering
    \setstretch{0.95}
    {\small\bfseries Министерство науки и высшего образования Российской Федерации} \\[2pt]
    {\small\bfseries Федеральное государственное автономное образовательное учреждение} \\[2pt]
    {\small\bfseries высшего образования} \\[2pt]
    {\small\bfseries «Московский государственный технический университет} \\[2pt]
    {\small\bfseries имени Н.Э. Баумана} \\[2pt]
    {\small\bfseries (национальный исследовательский университет)»} \\[2pt]
    {\small\bfseries (МГТУ им. Н.Э. Баумана)}
\end{minipage}

\vspace{12pt}
\hrule height 1pt
\vspace{2pt}
\hrule height 2pt

\vspace{10mm}

\noindent
\setstretch{1.2}
\textbf{ФАКУЛЬТЕТ} \underline{\makebox[138mm][l]{\hspace{15mm}\textbf{ИНФОРМАТИКА И СИСТЕМЫ УПРАВЛЕНИЯ}}} \\[10pt]
\textbf{КАФЕДРА} \underline{\makebox[143mm][l]{\hspace{12mm}\textbf{СИСТЕМЫ ОБРАБОТКИ ИНФОРМАЦИИ И УПРАВЛЕНИЯ}}}

\vspace{25mm}

\begin{center}
    \setstretch{1.3}
    {\Large Отчет по лабораторной работе № 4} \\[12pt]
    {\Large\bfseries «Линейные модели, SVM} \\[4pt]
    {\Large\bfseries и деревья решений»} \\[20mm]
    {\large по дисциплине «Технологии машинного обучения»}
\end{center}

\vspace{35mm}

\noindent
\setstretch{1.1}
\begin{tabular}{@{}llcc}
    \large Студент & \underline{\makebox[28mm][c]{\textbf{ИУ5-64Б}}} & \underline{\makebox[55mm]{}} & \underline{\makebox[50mm][c]{\textbf{Н.А. Чернев}}} \\[2pt]
    & \multicolumn{1}{c}{\scriptsize (Группа)} & \scriptsize (Подпись, дата) & \scriptsize (И.О. Фамилия) \\[20pt]
    \large Преподаватель & & \underline{\makebox[55mm]{}} & \underline{\makebox[50mm][c]{\textbf{Ю.Е. Гапанюк}}} \\[2pt]
    & & \scriptsize (Подпись, дата) & \scriptsize (И.О. Фамилия)
\end{tabular}

\vfill
\begin{center}
    \large 2026 г.
\end{center}

\newpage

\pagestyle{plain}

\section*{Цель работы}
Сравнение линейных моделей, SVM и деревьев решений на задаче классификации/регрессии.

\section*{Обучение моделей}

\subsection*{Логистическая регрессия}
\begin{lstlisting}
from sklearn.linear_model import LogisticRegression

log_reg = LogisticRegression(max_iter=1000, random_state=42)
log_reg.fit(X_train, y_train)
y_pred_lr = log_reg.predict(X_test)
\end{lstlisting}

\subsection*{Support Vector Machine (SVM)}
\begin{lstlisting}
from sklearn.svm import SVC

svm = SVC(kernel='rbf', C=1.0, gamma='scale', random_state=42)
svm.fit(X_train, y_train)
y_pred_svm = svm.predict(X_test)
\end{lstlisting}

\subsection*{Decision Tree}
\begin{lstlisting}
from sklearn.tree import DecisionTreeClassifier

dt = DecisionTreeClassifier(max_depth=5, random_state=42)
dt.fit(X_train, y_train)
y_pred_dt = dt.predict(X_test)
\end{lstlisting}

\section*{Оценка качества}

\begin{lstlisting}
from sklearn.metrics import accuracy_score, f1_score

models = {'LogReg': y_pred_lr, 'SVM': y_pred_svm, 'DT': y_pred_dt}

for name, y_pred in models.items():
    acc = accuracy_score(y_test, y_pred)
    f1 = f1_score(y_test, y_pred, average='weighted')
    print(f"{name}: Accuracy={acc:.4f}, F1={f1:.4f}")
\end{lstlisting}

\section*{Результаты сравнения}

\begin{center}
\begin{tabular}{|l|c|c|}
\hline
\textbf{Модель} & \textbf{Accuracy} & \textbf{F1-score} \\
\hline
Логистическая регрессия & 0.9667 & 0.9667 \\
\hline
SVM (RBF) & 0.9778 & 0.9778 \\
\hline
Decision Tree & 0.9667 & 0.9667 \\
\hline
\end{tabular}
\end{center}

\section*{Важность признаков и визуализация}

\begin{lstlisting}
import matplotlib.pyplot as plt
from sklearn.tree import plot_tree

# Важность признаков в дереве решений
importances = dt.feature_importances_
plt.figure(figsize=(10, 6))
plt.barh(range(len(importances)), importances)
plt.xlabel('Важность')
plt.ylabel('Признак')
plt.title('Важность признаков в Decision Tree')
plt.savefig('feature_importance.png', dpi=150, bbox_inches='tight')
plt.show()
\end{lstlisting}

\begin{figure}[h!]
    \centering
    \includegraphics[width=0.85\textwidth]{fig_1.png}
    \newline
    \small Рисунок 1 --- Важность признаков в дереве решений
    \newline
\end{figure}

\subsection*{Визуализация структуры дерева}
\begin{lstlisting}
# Визуализация структуры дерева
plt.figure(figsize=(20, 10))
plot_tree(dt, feature_names=feature_names, class_names=class_names, 
          filled=True, rounded=True)
plt.title('Структура Decision Tree')
plt.tight_layout()
plt.savefig('tree_structure.png', dpi=150, bbox_inches='tight')
plt.show()
\end{lstlisting}

\begin{figure}[h!]
    \centering
    \includegraphics[width=0.85\textwidth]{fig_2.png}
    \newline
    \small Рисунок 2 --- Структура дерева решений на датасете Iris
    \newline
\end{figure}

\section*{Выводы}

\begin{enumerate}
    \item SVM показала лучшую точность (0.9778) благодаря нелинейному ядру RBF
    \item Decision Tree обеспечивает интерпретируемые правила решения
    \item Логистическая регрессия быстрее, но менее гибкая
    \item Для данного датасета (Iris) все три модели работают хорошо
\end{enumerate}

\end{document}
